\documentclass[11pt]{article}
\usepackage[a4paper,margin=1in]{geometry}
\usepackage{enumitem}
\usepackage{hyperref}
\usepackage{array}
\usepackage{longtable}
\usepackage{titlesec}

\titleformat{\section}{\large\bfseries}{}{0em}{}
\titleformat{\subsection}{\normalsize\bfseries}{}{0em}{}

\setlist[itemize]{leftmargin=1.4em}
\setlist[enumerate]{leftmargin=1.6em}

\title{Presentation Outline: DeepRL Self-Play Arena}
\author{Julie Yang}
\date{\today}

\begin{document}
\maketitle

\section{Goal of This Outline}

This document is a slide-by-slide outline for the final presentation of the DeepRL Self-Play Arena project. It is not a final report. The emphasis is therefore on:

\begin{itemize}
    \item motivation and project scope,
    \item development process and engineering milestones,
    \item major design decisions,
    \item demonstrations and experimental observations,
    \item what worked, what did not work, and what was learned.
\end{itemize}

The presentation should feel like a story of building a working multi-agent RL system in Unity, not just a list of algorithms.

\section{Recommended High-Level Presentation Structure}

Recommended total length: 12--16 slides.

\begin{enumerate}
    \item Motivation and project definition
    \item Related work / inspiration
    \item Environment and task design
    \item System architecture
    \item Development milestones
    \item Training pipeline evolution
    \item Results and demos
    \item Limitations and lessons learned
    \item Final takeaways
\end{enumerate}

\section{Detailed Slide Outline}

\subsection*{Slide 1: Title Slide}

\textbf{Title:} DeepRL Self-Play Arena: Building a Unity-Based Multi-Agent RL Environment

\textbf{Main points}
\begin{itemize}
    \item Course project title
    \item Your name
    \item Course / semester
    \item One-line summary: a 1v1 sumo-style self-play RL environment with training in Python and real-time demonstration in Unity
\end{itemize}

\textbf{Suggested visual}
\begin{itemize}
    \item Screenshot of the Unity arena
    \item Or a screenshot of human-vs-policy mode
\end{itemize}

\textbf{Speaker goal}
\begin{itemize}
    \item Give the audience immediate intuition for the project before any technical explanation
\end{itemize}

\subsection*{Slide 2: Motivation}

\textbf{Suggested title:} Why This Project?

\textbf{Main points}
\begin{itemize}
    \item Self-play is powerful but unstable in competitive multi-agent RL
    \item Many classic examples, such as AlphaGo Zero and OpenAI Five, show the strength of self-play but also imply non-stationarity and opponent-distribution issues
    \item A small Unity arena is a good testbed:
    \begin{itemize}
        \item simple enough to debug,
        \item rich enough to exhibit policy collapse, cycling, and brittle behaviors,
        \item visually interpretable for presentation
    \end{itemize}
    \item Personal motivation: build not only a training script, but a complete interactive system that supports both learning and demonstration
\end{itemize}

\textbf{Suggested phrasing}
\begin{itemize}
    \item The project asks: can we build a compact competitive RL environment that is easy to visualize, yet still exposes the core stability problems of self-play?
\end{itemize}

\subsection*{Slide 3: Original Proposal vs. Final Development Direction}

\textbf{Suggested title:} Initial Plan and How It Changed

\textbf{Main points}
\begin{itemize}
    \item Original proposal focused on:
    \begin{itemize}
        \item Unity ML-Agents style PPO self-play
        \item opponent pooling and sampling strategies
        \item evaluating training stability and robustness
    \end{itemize}
    \item In practice, the early challenge was not only algorithm design but full-stack integration
    \item The project direction shifted toward:
    \begin{itemize}
        \item building a reliable Unity--Python bridge,
        \item defining a stable observation/action interface,
        \item constructing end-to-end training and evaluation loops,
        \item then iterating on self-play structure and reward design
    \end{itemize}
    \item This was a necessary shift, not a distraction
\end{itemize}

\textbf{Key message}
\begin{itemize}
    \item The biggest real project lesson was that systems integration dominated algorithm experimentation in the first half of development
\end{itemize}

\subsection*{Slide 4: Literature and Inspiration}

\textbf{Suggested title:} Related Work and What This Project Borrowed

\textbf{Main points}
\begin{itemize}
    \item AlphaGo Zero:
    \begin{itemize}
        \item self-play as a mechanism for continual policy improvement
    \end{itemize}
    \item OpenAI Five:
    \begin{itemize}
        \item importance of training against historical versions and maintaining opponent diversity
    \end{itemize}
    \item Emergent Complexity via Multi-Agent Competition:
    \begin{itemize}
        \item multi-agent competition can create interesting behaviors, but also unstable or degenerate training dynamics
    \end{itemize}
    \item This project did not attempt to reproduce those systems
    \item Instead, it translated their core ideas into a much smaller and more inspectable Unity environment
\end{itemize}

\textbf{Suggested visual}
\begin{itemize}
    \item A simple table with three rows:
    \begin{itemize}
        \item paper/system,
        \item idea borrowed,
        \item how it influenced this project
    \end{itemize}
\end{itemize}

\subsection*{Slide 5: Task and Environment Definition}

\textbf{Suggested title:} The Arena Task

\textbf{Main points}
\begin{itemize}
    \item Two circular agents compete inside a circular arena
    \item Objective: push the opponent outside the arena boundary
    \item Episode ends on:
    \begin{itemize}
        \item opponent ring-out,
        \item self ring-out,
        \item time limit
    \end{itemize}
    \item Agent action semantics:
    \begin{itemize}
        \item 2D movement,
        \item push/dash with cooldown
    \end{itemize}
    \item Why this task is useful:
    \begin{itemize}
        \item easy to understand,
        \item compact state/action space,
        \item still contains positioning, pressure, risk control, and timing
    \end{itemize}
\end{itemize}

\textbf{Suggested visual}
\begin{itemize}
    \item top-down arena screenshot with labels for spawn points, arena boundary, and agents
\end{itemize}

\subsection*{Slide 6: System Architecture}

\textbf{Suggested title:} Unity + Python Architecture

\textbf{Main points}
\begin{itemize}
    \item Unity handles:
    \begin{itemize}
        \item physics,
        \item arena state,
        \item reset logic,
        \item final demo presentation
    \end{itemize}
    \item Python handles:
    \begin{itemize}
        \item rollout collection,
        \item observation/action vectorization,
        \item imitation learning,
        \item PPO training,
        \item checkpoint evaluation
    \end{itemize}
    \item The two are connected by a custom TCP bridge
\end{itemize}

\textbf{Suggested subdiagram}
\begin{itemize}
    \item Unity scene $\leftrightarrow$ TCP bridge $\leftrightarrow$ Python env wrapper $\leftrightarrow$ training / evaluation scripts
\end{itemize}

\textbf{Key message}
\begin{itemize}
    \item One of the central achievements of the project is not a single policy result, but a working interactive RL stack
\end{itemize}

\subsection*{Slide 7: Development Phase 1 --- Bridging Python and Unity}

\textbf{Suggested title:} Milestone 1: Making Unity and Python Talk

\textbf{Main points}
\begin{itemize}
    \item Built a TCP bridge for command/response interaction
    \item Early milestone:
    \begin{itemize}
        \item send actions from Python,
        \item advance the Unity environment,
        \item return observations, rewards, and done flags
    \end{itemize}
    \item Key debugging issues:
    \begin{itemize}
        \item connection stability,
        \item request/response format,
        \item persistent vs. short-lived connections,
        \item distinguishing training mode from real-time mode
    \end{itemize}
    \item Final result:
    \begin{itemize}
        \item persistent TCP connection,
        \item single-environment control,
        \item later extension to batched arenas
    \end{itemize}
\end{itemize}

\textbf{Suggested visual}
\begin{itemize}
    \item simple request/response JSON example
\end{itemize}

\subsection*{Slide 8: Development Phase 2 --- Action Format}

\textbf{Suggested title:} Milestone 2: Defining the Action Space

\textbf{Main points}
\begin{itemize}
    \item Initial task: define actions that are simple enough for RL but still expressive
    \item Final vector action:
    \begin{itemize}
        \item `move_x`
        \item `move_y`
        \item `use_push`
    \end{itemize}
    \item Important design choice:
    \begin{itemize}
        \item push direction depends on current motion direction rather than a separate push vector
    \end{itemize}
    \item Why this was useful:
    \begin{itemize}
        \item smaller action space,
        \item easier policy output semantics,
        \item more stable interface between Python and Unity
    \end{itemize}
    \item Later refinement:
    \begin{itemize}
        \item move became continuous/held input,
        \item push remained pulse-like
    \end{itemize}
\end{itemize}

\textbf{Nice insight to say aloud}
\begin{itemize}
    \item A lot of RL difficulty comes from badly chosen action semantics, not just from the optimizer
\end{itemize}

\subsection*{Slide 9: Development Phase 3 --- Observation Design}

\textbf{Suggested title:} Milestone 3: What the Agent Sees

\textbf{Main points}
\begin{itemize}
    \item Observation started simple:
    \begin{itemize}
        \item self position / velocity
        \item opponent position / velocity
        \item push readiness
    \end{itemize}
    \item Then expanded from 13D to 17D
    \item Added explicit arena-awareness:
    \begin{itemize}
        \item self distance to center
        \item self edge margin
        \item opponent distance to center
        \item opponent edge margin
    \end{itemize}
    \item Later added:
    \begin{itemize}
        \item left-right canonicalization,
        \item opponent-position scaling,
        \item relative-position scaling
    \end{itemize}
    \item Why this mattered:
    \begin{itemize}
        \item reduced symmetry-breaking artifacts,
        \item reduced strange side-dependent behaviors,
        \item made the opponent's location more salient
    \end{itemize}
\end{itemize}

\textbf{Suggested visual}
\begin{itemize}
    \item one slide with the 17D vector grouped into categories
\end{itemize}

\subsection*{Slide 10: Development Phase 4 --- Data and Evaluation Pipeline}

\textbf{Suggested title:} Milestone 4: Rollouts, Checkpoints, and Evaluation

\textbf{Main points}
\begin{itemize}
    \item Built a formal Python environment wrapper:
    \begin{itemize}
        \item `reset`
        \item `step`
        \item `close`
    \end{itemize}
    \item Added rollout collection
    \item Added JSON trajectory saving
    \item Added checkpoint evaluation scripts
    \item Added multi-episode evaluation summaries
    \item This made the project experimentally usable rather than only visually interactive
\end{itemize}

\textbf{Key message}
\begin{itemize}
    \item Before PPO experiments mattered, the project needed reliable instrumentation
\end{itemize}

\subsection*{Slide 11: Development Phase 5 --- Imitation Learning Baseline}

\textbf{Suggested title:} Milestone 5: Sanity Check via Behavior Cloning

\textbf{Main points}
\begin{itemize}
    \item Before full RL, trained a small MLP policy on heuristic rollouts
    \item Why imitation learning was useful:
    \begin{itemize}
        \item checked observation formatting,
        \item checked action decoding,
        \item checked checkpoint save/load,
        \item checked evaluation back in Unity
    \end{itemize}
    \item Result:
    \begin{itemize}
        \item the learned MLP could reliably imitate simple heuristic behavior
    \end{itemize}
\end{itemize}

\textbf{Interpretation}
\begin{itemize}
    \item This did not solve the final task, but it validated the core learning pipeline
\end{itemize}

\subsection*{Slide 12: Development Phase 6 --- PPO Baseline}

\textbf{Suggested title:} Milestone 6: From Imitation to PPO

\textbf{Main points}
\begin{itemize}
    \item Added a minimal actor-critic PPO implementation
    \item First tested in a simpler setting:
    \begin{itemize}
        \item train one agent against a fixed scripted opponent
    \end{itemize}
    \item This helped verify:
    \begin{itemize}
        \item rollout collection under on-policy learning,
        \item value estimation,
        \item clipping loss,
        \item checkpoint persistence
    \end{itemize}
    \item Observed issue:
    \begin{itemize}
        \item training could run, but learned behavior was often weak or degenerate
    \end{itemize}
\end{itemize}

\textbf{Suggested plot}
\begin{itemize}
    \item one example of PPO log values or a screenshot of console outputs with policy loss / value loss / returns
\end{itemize}

\subsection*{Slide 13: Development Phase 7 --- Self-Play Variants}

\textbf{Suggested title:} Milestone 7: Shared Self-Play, Then Alternating A/B

\textbf{Main points}
\begin{itemize}
    \item Started with shared-parameter self-play
    \item Problem:
    \begin{itemize}
        \item too easy to collapse into symmetric low-quality behavior
    \end{itemize}
    \item Upgraded to alternating two-policy training:
    \begin{itemize}
        \item policy A trains while B is frozen,
        \item then policy B trains while A is frozen
    \end{itemize}
    \item Additional mechanisms introduced later:
    \begin{itemize}
        \item finish active episodes before segment exit,
        \item EMA propagation between A and B,
        \item opponent pool sampling,
        \item periodic mutation
    \end{itemize}
\end{itemize}

\textbf{Key message}
\begin{itemize}
    \item A large part of the project became designing a better self-play training process, not just a better neural network
\end{itemize}

\subsection*{Slide 14: Development Phase 8 --- Batch PPO}

\textbf{Suggested title:} Milestone 8: Parallel Arenas for Higher Throughput

\textbf{Main points}
\begin{itemize}
    \item Single-arena training was too slow and too correlated
    \item Added a batched 4-arena Unity mode
    \item Key engineering issue:
    \begin{itemize}
        \item batch stepping must share one global physics simulation and advance all arenas together
    \end{itemize}
    \item Added seeded reset with per-arena variation
    \item Later increased arena count in config
    \item Goal:
    \begin{itemize}
        \item more data throughput,
        \item less repeated experience,
        \item better diversity during training
    \end{itemize}
\end{itemize}

\textbf{Suggested visual}
\begin{itemize}
    \item screenshot of multiple arenas running at once
\end{itemize}

\subsection*{Slide 15: Development Phase 9 --- Human vs Policy Demo}

\textbf{Suggested title:} Milestone 9: Real-Time Play Against the Learned Policy

\textbf{Main points}
\begin{itemize}
    \item Added a separate real-time mode:
    \begin{itemize}
        \item human controls agent 0 with keyboard,
        \item Python drives agent 1 using the latest checkpoint
    \end{itemize}
    \item This required separating:
    \begin{itemize}
        \item training-time manual simulation mode,
        \item real-time interactive mode
    \end{itemize}
    \item Why this mattered:
    \begin{itemize}
        \item better qualitative evaluation,
        \item much more compelling final demo,
        \item exposed weaknesses that pure metrics did not show
    \end{itemize}
\end{itemize}

\textbf{This should definitely be a live demo or a short video}

\subsection*{Slide 16: Reward Design Evolution}

\textbf{Suggested title:} Reward Design Was Harder Than Expected

\textbf{Main points}
\begin{itemize}
    \item Sparse win/loss reward is clean but weak
    \item Heavy draw or timeout penalties caused bad side effects
    \begin{itemize}
        \item too passive,
        \item too suicidal,
        \item strange edge-seeking behavior
    \end{itemize}
    \item Introduced several refinements during development:
    \begin{itemize}
        \item edge safety shaping,
        \item outward pressure shaping,
        \item time-scaled loss penalty,
        \item timeout center-bias penalty
    \end{itemize}
    \item Main lesson:
    \begin{itemize}
        \item the difficult part was not only defining the objective ``win,'' but defining what useful intermediate behavior should be rewarded
    \end{itemize}
\end{itemize}

\textbf{Good thing to say out loud}
\begin{itemize}
    \item This project made it very clear that reward design is not a small detail; it shapes the entire character of the learned policy
\end{itemize}

\subsection*{Slide 17: Results --- What Was Successfully Achieved}

\textbf{Suggested title:} Main Results and Deliverables

\textbf{Main points}
\begin{itemize}
    \item Built a functioning Unity competitive environment
    \item Built a custom Unity--Python TCP integration layer
    \item Built a full training/evaluation workflow in Python
    \item Verified imitation learning and PPO checkpoint loops
    \item Implemented alternating two-policy self-play
    \item Implemented batched training arenas
    \item Implemented human-vs-policy real-time play mode
\end{itemize}

\textbf{Good framing}
\begin{itemize}
    \item Even when the final learned policy is imperfect, the system-level deliverable is substantial and functional
\end{itemize}

\subsection*{Slide 18: Results --- Example Behaviors and Qualitative Findings}

\textbf{Suggested title:} What the Agents Learned, and What They Still Struggle With

\textbf{Possible talking points}
\begin{itemize}
    \item Learned some basic pressure and arena-awareness behaviors
    \item Different reward settings led to clearly different styles
    \item Human-vs-policy mode made weaknesses obvious:
    \begin{itemize}
        \item brittle tactics,
        \item symmetry artifacts,
        \item overcommitment or passivity,
        \item weak adaptation against a real human
    \end{itemize}
    \item More complex self-play machinery improved diversity, but did not fully solve strategic sophistication
\end{itemize}

\textbf{Suggested visual}
\begin{itemize}
    \item GIF/video stills of:
    \begin{itemize}
        \item a successful push-out,
        \item a weird degenerate policy,
        \item human-vs-policy interaction
    \end{itemize}
\end{itemize}

\subsection*{Slide 19: What Was Not Fully Achieved}

\textbf{Suggested title:} Limitations

\textbf{Main points}
\begin{itemize}
    \item Final policies are still not deeply strategic against humans
    \item Reward design remains partially heuristic
    \item Opponent pool and mutation are implemented, but not yet exhaustively studied
    \item The project did not finish with a clean large-scale ablation over all self-play variants
    \item This is still closer to a strong prototype and research platform than to a polished benchmark paper
\end{itemize}

\textbf{Important tone}
\begin{itemize}
    \item This slide should sound honest but not apologetic
    \item Emphasize that the key contribution is building and validating the platform, then exploring several nontrivial self-play design choices
\end{itemize}

\subsection*{Slide 20: Lessons Learned}

\textbf{Suggested title:} Main Takeaways

\textbf{Main points}
\begin{itemize}
    \item Multi-agent RL is as much a systems problem as an algorithm problem
    \item Observation design strongly affects whether training can even begin to stabilize
    \item Reward design is extremely sensitive in competitive environments
    \item Symmetry and opponent diversity matter a lot
    \item Visual debugging in Unity was genuinely useful, not just cosmetic
\end{itemize}

\textbf{Nice closing sentence}
\begin{itemize}
    \item The project ended up showing that even a very small arena can surface many of the same stability issues seen in much larger self-play systems
\end{itemize}

\subsection*{Slide 21: Closing / Demo Slide}

\textbf{Suggested title:} Final Demo and Questions

\textbf{Options}
\begin{itemize}
    \item End with a short live human-vs-policy demo
    \item Or end with a short recorded clip if live demo is risky
    \item Show the final Unity scene and briefly restate the project contributions
\end{itemize}

\textbf{Closing bullets}
\begin{itemize}
    \item Built a Unity-based competitive RL environment
    \item Implemented end-to-end learning and evaluation pipeline
    \item Explored multiple self-play stabilization mechanisms
    \item Produced an interactive final demonstration
\end{itemize}

\section{Recommended Slide Compression If Time Is Short}

If the presentation must be shorter, compress to roughly 12 slides using this mapping:

\begin{enumerate}
    \item Title
    \item Motivation + related work
    \item Environment/task
    \item System architecture
    \item Bridge and interface milestones
    \item Imitation + PPO pipeline
    \item Self-play evolution
    \item Batch training + human-vs-policy
    \item Reward design challenges
    \item Main results
    \item Limitations / lessons learned
    \item Demo / questions
\end{enumerate}

\section{What to Emphasize During Presentation}

\begin{itemize}
    \item Do not oversell final policy strength.
    \item Do emphasize how much engineering was required to turn the idea into a working system.
    \item Use the development chronology as the main narrative.
    \item Make Unity screenshots, bridge diagrams, and human-vs-policy visuals do real explanatory work.
    \item When discussing results, separate clearly:
    \begin{itemize}
        \item what is fully implemented,
        \item what was experimentally promising,
        \item what remains incomplete.
    \end{itemize}
\end{itemize}

\section{Suggested Assets to Prepare for the Actual Slides}

\begin{itemize}
    \item One clean arena screenshot
    \item One architecture diagram
    \item One bridge JSON or command flow illustration
    \item One observation/action slide graphic
    \item One screenshot of multi-arena batched training
    \item One screenshot of human-vs-policy mode
    \item One or two console/log screenshots showing PPO updates
    \item One short GIF or video clip of the final demo
\end{itemize}

\section{Suggested Final Narrative in One Sentence}

This project started as a study of self-play stability, and evolved into the construction of a full Unity--Python competitive RL system that supports environment design, self-play experimentation, batched training, and real-time human-vs-policy demonstration.

\end{document}
